% ============================================================================
% CHAPTER: PROPOSED ARCHITECTURE AND MASK-BASED DENSE ALIGNMENT
% ============================================================================
% This chapter is intended to be \input{} or \include{} from a main document
% that defines \documentclass, \usepackage, and \begin{document}.
% Required packages: amsmath, amssymb, graphicx, booktabs, algorithm2e or
% algorithmic, hyperref.  A \bibliography{references} call should appear in
% the main document.
% ============================================================================

\chapter{Proposed Architecture and Mask-Based Dense Alignment}
\label{ch:proposed_architecture}

\section{Introduction and Design Rationale}
\label{sec:arch_intro}

Classical approaches to infrared--visible image registration rely on sparse
keypoint detection followed by robust estimation of a global geometric
transformation---typically an affine map or projective homography
\cite{zitova2003image}.  While effective when the scene is approximately
planar or when the two sensors share a narrow baseline, such methods fail in
precisely the situations that motivate multi-spectral fusion: parallax from
objects at varying depths, independently moving targets, and the
radiometric dissimilarity between long-wave infrared (LWIR) and visible
(RGB) imagery that causes feature detectors to return inconsistent
keypoints \cite{brown2007automatic}.

This chapter develops a \emph{dense, learned alignment} framework that
predicts a per-pixel two-dimensional displacement field
$\mathbf{d}: \Omega \rightarrow \mathbb{R}^{2}$, where $\Omega \subset
\mathbb{Z}^{2}$ is the discrete image lattice of spatial dimensions
$H \times W$.  Given an infrared frame $I_{\mathrm{IR}} \in
\mathbb{R}^{H \times W \times 1}$ and a visible frame $I_{\mathrm{RGB}}
\in \mathbb{R}^{H \times W \times 3}$, the network produces a field
$\hat{\mathbf{d}} \in \mathbb{R}^{H \times W \times 2}$ such that
warping $I_{\mathrm{RGB}}$ by $\hat{\mathbf{d}}$ yields maximal spatial
agreement with $I_{\mathrm{IR}}$.  Crucially, the network also emits a
\emph{confidence mask} $\hat{m} \in [0,1]^{H \times W}$ that indicates
which displacement predictions are reliable.  This mask-based mechanism
allows the system to gracefully handle occlusions, spectral ambiguities,
and textureless regions where no unique correspondence exists.

The architecture draws inspiration from the dense optical-flow literature,
particularly RAFT \cite{teed2020raft} and PWC-Net \cite{sun2018pwcnet},
but introduces domain-specific modifications: a dual-encoder front end that
respects the distinct statistical properties of IR and RGB imagery, a
learnable confidence head, and a semantic-mask gating module for
object-aware alignment.

% ========================================================================
\section{Overall Architecture Design}
\label{sec:overall_arch}

% ------------------------------------------------------------------------
\subsection{Dual-Encoder Feature Extraction}
\label{subsec:dual_encoder}

Because the intensity distributions of thermal-infrared and visible-light
images differ fundamentally---IR encodes emissivity and temperature while
RGB encodes reflectance---sharing early convolutional filters between the
two modalities is suboptimal \cite{li2018densefuse}.  We therefore adopt a
\textbf{dual-encoder} topology in which each modality is processed by its
own feature backbone before the two representations are merged.

Each encoder $\mathcal{E}_{\mathrm{IR}}$ and $\mathcal{E}_{\mathrm{RGB}}$
is instantiated as a ResNet-18 \cite{he2016deep} truncated after the
fourth residual stage, producing a feature pyramid with outputs at
$1/4$, $1/8$, $1/16$, and $1/32$ of the input spatial resolution.
Concretely, for an input of size $H \times W$ (we use $H{=}W{=}512$ in
all experiments), the feature maps at each level $l \in \{1,2,3,4\}$ have
spatial dimensions $H/2^{l+1} \times W/2^{l+1}$ and channel counts
$C_l \in \{64,128,256,512\}$ respectively.  Each backbone is initialised
from ImageNet-pretrained weights; the single-channel IR input is
accommodated by averaging the three channels of the pretrained first
convolutional kernel into one and replicating the result, following the
protocol of \cite{zhao2023cddfuse}.

An \textbf{alternative single-encoder design} concatenates the IR and RGB
images along the channel axis to form a four-channel input
$[I_{\mathrm{IR}}; I_{\mathrm{RGB}}] \in \mathbb{R}^{H \times W \times
4}$, which is then processed by a U-Net-style encoder--decoder
\cite{ronneberger2015unet}.  This approach has fewer parameters and lower
latency, but our ablation study (Chapter~\ref{ch:experiments}) shows that
it consistently underperforms the dual-encoder variant by 0.4--0.7\,px in
endpoint error, confirming that modality-specific feature learning is
beneficial.

% ------------------------------------------------------------------------
\subsection{Dense Displacement Field Output}
\label{subsec:displacement_field}

The decoder produces a displacement field $\hat{\mathbf{d}} \in
\mathbb{R}^{H \times W \times 2}$, where each spatial location
$(u,v)$ stores a horizontal shift $d_x(u,v)$ and a vertical shift
$d_y(u,v)$.  The aligned RGB image is obtained by a
\textbf{differentiable bilinear sampler}:
\begin{equation}
  \hat{I}_{\mathrm{RGB}}^{\,\text{warp}}(u,v)
    = I_{\mathrm{RGB}}\!\bigl(u + d_x(u,v),\; v + d_y(u,v)\bigr),
  \label{eq:warp}
\end{equation}
implemented via \texttt{torch.nn.functional.grid\_sample} in PyTorch
\cite{jaderberg2015spatial}.  Because \texttt{grid\_sample} is
differentiable with respect to both the input image and the sampling grid,
the entire pipeline is end-to-end trainable.  Displacement values are
predicted in pixel units and converted to the $[-1,1]$ normalised
coordinate system required by the API at inference time.

% ========================================================================
\section{Mask-Guided Alignment}
\label{sec:mask_guided}

% ------------------------------------------------------------------------
\subsection{Confidence Mask Prediction}
\label{subsec:confidence_mask}

Not every pixel admits a unique, well-defined correspondence between the
two modalities.  Specular highlights visible only in RGB, thermal blooming
in IR, and occluded regions arising from parallax all produce locations
where the ground-truth displacement is either undefined or inherently
ambiguous.  Rather than forcing the network to hallucinate displacements
for such pixels, we augment the output with a \textbf{confidence mask}
$\hat{m} \in [0,1]^{H \times W}$ produced by a parallel $1 \times 1$
convolutional head followed by a sigmoid activation.

The combined output is therefore $\mathbb{R}^{H \times W \times 3}$: two
channels for the displacement $(d_x, d_y)$ and one channel for the
confidence $m$.  During training, the photometric and geometric losses
(Chapter~\ref{ch:training}) are weighted by $\hat{m}$:
\begin{equation}
  \mathcal{L}_{\text{align}}
    = \frac{1}{|\Omega|} \sum_{(u,v) \in \Omega}
      \hat{m}(u,v) \cdot
      \bigl\| \hat{\mathbf{d}}(u,v) - \mathbf{d}^{\ast}(u,v) \bigr\|_1
      \;-\; \lambda \log \hat{m}(u,v),
  \label{eq:mask_loss}
\end{equation}
where $\mathbf{d}^{\ast}$ is the ground-truth displacement and $\lambda$
is a regularisation coefficient (set to $0.5$ in our experiments).  The
first term encourages the network to suppress confident predictions at
unreliable locations; the second term---a log-barrier---prevents the
trivial solution $\hat{m} \equiv 0$.  This formulation follows the
self-supervised uncertainty framework of Kendall and Gal
\cite{kendall2017uncertainties}, adapted here to a registration context.

% ------------------------------------------------------------------------
\subsection{Semantic-Mask-Guided Alignment}
\label{subsec:semantic_mask}

When pre-computed semantic segmentation masks are available (e.g.\ from a
separately trained Mask R-CNN \cite{he2017maskrcnn} or
SAM~\cite{kirillov2023sam}), alignment accuracy can be improved by
treating each semantic region independently.  We define a set of binary
masks $\{S_k\}_{k=1}^{K}$ corresponding to $K$ semantic classes (e.g.\
person, vehicle, building, sky) and compute per-class alignment losses:
\begin{equation}
  \mathcal{L}_{\text{sem}}
    = \sum_{k=1}^{K} w_k \;
      \frac{\sum_{(u,v)} S_k(u,v) \cdot
        \|\hat{\mathbf{d}}(u,v) - \mathbf{d}^{\ast}(u,v)\|_1}
       {\sum_{(u,v)} S_k(u,v) + \epsilon},
  \label{eq:sem_loss}
\end{equation}
with class weights $w_k$ set inversely proportional to class area.  This
scheme ensures that small but safety-critical objects (pedestrians,
signage) receive alignment supervision commensurate with their importance,
rather than being dominated by large background regions.

Additionally, the semantic masks are injected into the decoder as a
\textbf{gating signal}.  After the last skip connection in the decoder,
the concatenated features are element-wise multiplied by a learned
projection of the one-hot semantic map, producing class-conditioned
displacement predictions.  This allows the network to learn, for example,
that sky regions tolerate large uncertainty while object boundaries demand
sub-pixel precision.

% ------------------------------------------------------------------------
\subsection{Spatial Attention for Salient-Region Focus}
\label{subsec:attention}

Even without explicit semantic labels, the network benefits from focusing
its representational capacity on informative regions.  We insert a
lightweight \textbf{convolutional block attention module}
(CBAM)~\cite{woo2018cbam} after each encoder stage.  CBAM produces a
channel attention vector via global average- and max-pooling followed by a
shared MLP, and a spatial attention map via a $7 \times 7$ convolution
over the channel-pooled feature map.  The resulting spatial attention map
up-weights edges, object contours, and thermal gradients while
down-weighting uniform regions (sky, road surface), where alignment cues
are weakest.  In our ablations (Section~\ref{sec:ablation}), adding CBAM
reduces mean endpoint error by 8\% relative to the base dual-encoder
without attention.

% ========================================================================
\section{Correlation-Based Feature Matching}
\label{sec:correlation}

% ------------------------------------------------------------------------
\subsection{4D Correlation Volume Construction}
\label{subsec:corr_volume}

Given feature maps $\mathbf{F}_{\mathrm{IR}} \in \mathbb{R}^{H' \times
W' \times D}$ and $\mathbf{F}_{\mathrm{RGB}} \in \mathbb{R}^{H' \times
W' \times D}$ extracted at the $1/8$ resolution level ($H' = H/8$,
$W' = W/8$, $D = 128$ after a $1 \times 1$ projection), we construct a
\textbf{4D correlation volume}:
\begin{equation}
  \mathbf{C}(u_1, v_1, u_2, v_2)
    = \frac{
        \mathbf{F}_{\mathrm{IR}}(u_1, v_1)^{\!\top}
        \mathbf{F}_{\mathrm{RGB}}(u_2, v_2)
      }{D},
  \label{eq:corr_volume}
\end{equation}
yielding a tensor $\mathbf{C} \in \mathbb{R}^{H' \times W' \times H'
\times W'}$.  For $H'{=}W'{=}64$ this requires $64^4 \approx 16.8$\,M
entries, which is tractable at single precision (roughly 67\,MB).
Following RAFT~\cite{teed2020raft}, we construct a \emph{correlation
pyramid} by average-pooling the last two dimensions of $\mathbf{C}$ with
kernel sizes $\{1, 2, 4, 8\}$, yielding four volumes at progressively
coarser matching resolution.  At query time the network looks up a local
neighbourhood in each pyramid level centred on the current displacement
estimate, concatenates the results, and feeds them into the update
operator.

The cross-modal nature of IR--RGB matching makes the raw dot-product
correlation noisier than in the mono-spectral optical-flow setting.  To
mitigate this, we apply a two-layer MLP with ReLU activation to each
feature vector before computing the inner product, following the
\emph{learned matching cost} strategy of \cite{xu2017accurate}.

% ------------------------------------------------------------------------
\subsection{Iterative Refinement via Recurrent Update}
\label{subsec:iterative_refinement}

Rather than regressing the full displacement field in a single forward
pass, we adopt the iterative refinement paradigm of RAFT
\cite{teed2020raft}.  A convolutional gated recurrent unit (GRU)
maintains a hidden state $\mathbf{h}_t \in \mathbb{R}^{H' \times W'
\times 128}$ and at each iteration $t = 1, \ldots, T$ produces a
displacement update $\Delta \mathbf{d}_t$:
\begin{equation}
  \hat{\mathbf{d}}_{t} = \hat{\mathbf{d}}_{t-1} + \Delta \mathbf{d}_{t},
  \qquad
  \Delta \mathbf{d}_{t} = \text{GRU}\!\bigl(
    \mathbf{h}_{t-1},\;
    [\,\text{corr}(\hat{\mathbf{d}}_{t-1});\;
     \hat{\mathbf{d}}_{t-1};\;
     \mathbf{c}\,]
  \bigr),
  \label{eq:gru_update}
\end{equation}
where $\text{corr}(\hat{\mathbf{d}}_{t-1})$ denotes the correlation
look-up indexed by the current displacement estimate and $\mathbf{c}$
is a context feature map extracted from $I_{\mathrm{IR}}$.  The GRU is
unrolled for $T{=}12$ iterations during training and $T{=}24$ at
inference; each iteration adds negligible memory because the correlation
volume is pre-computed and only indexed, not recomputed.

Each intermediate estimate $\hat{\mathbf{d}}_t$ is supervised with
exponentially increasing weight $\gamma^{T-t}$ ($\gamma = 0.85$),
encouraging the network to converge quickly:
\begin{equation}
  \mathcal{L}_{\text{iter}}
    = \sum_{t=1}^{T} \gamma^{T-t}\,
      \|\hat{\mathbf{d}}_t - \mathbf{d}^{\ast}\|_1.
  \label{eq:iter_loss}
\end{equation}

% ========================================================================
\section{Multi-Scale and Coarse-to-Fine Strategy}
\label{sec:multiscale}

% ------------------------------------------------------------------------
\subsection{Motivation}
\label{subsec:multiscale_motivation}

Real-world IR--RGB misalignment spans a wide range of magnitudes.
Mechanical vibration or thermal expansion of the sensor mount can
introduce global shifts of 20--50 pixels, while parallax-induced
displacement near depth discontinuities varies smoothly within a
5--10\,pixel range.  Capturing both regimes within a single network
requires receptive fields that are simultaneously large (to detect coarse
global offset) and precise (to recover sub-pixel local detail).

A coarse-to-fine architecture addresses this by decomposing the estimation
problem across spatial scales.  At the coarsest level ($1/32$ resolution,
$16 \times 16$ for a $512 \times 512$ input) the effective receptive field
of a $3 \times 3$ convolution stack covers the entire image, enabling
detection of large translations and rotations.  At the finest level
($1/4$ resolution, $128 \times 128$) the same convolutional stack resolves
sub-pixel displacements around edges and fine structures.

% ------------------------------------------------------------------------
\subsection{Feature Pyramid Network}
\label{subsec:fpn}

We augment each ResNet-18 encoder with a lightweight Feature Pyramid
Network (FPN) \cite{lin2017fpn}.  Top-down lateral connections fuse
high-level semantic features with low-level spatial detail, producing
pyramid levels $\{\mathbf{P}_l\}_{l=1}^{4}$ each with a uniform channel
dimension of $D{=}128$.  The FPN adds less than 1\,M parameters per
encoder and introduces negligible latency.

% ------------------------------------------------------------------------
\subsection{Coarse-to-Fine Displacement Estimation}
\label{subsec:coarse_to_fine}

Displacement estimation proceeds from level $l{=}4$ (coarsest) to
$l{=}1$ (finest).  At each level:
\begin{enumerate}
  \item The correlation volume is computed between
        $\mathbf{P}_l^{\mathrm{IR}}$ and $\mathbf{P}_l^{\mathrm{RGB}}$.
  \item The iterative GRU refiner runs for $T_l$ iterations
        ($T_4{=}6$, $T_3{=}4$, $T_2{=}3$, $T_1{=}3$).
  \item The resulting displacement field is bilinearly upsampled by
        $2\times$ and passed as initialisation to level $l{-}1$.
\end{enumerate}
This telescoping strategy halves the number of refinement iterations
required at fine scales (since the initialisation is already close to the
true displacement), reducing total compute by roughly 40\% relative to
running 12 iterations at a single scale.

% ========================================================================
\section{Concrete Architecture Specification}
\label{sec:arch_spec}

Table~\ref{tab:architecture} provides a layer-by-layer specification of
the complete network, which we refer to as \textbf{IRAlign-Net}.

\begin{table}[t]
\centering
\caption{IRAlign-Net architecture summary for $512 \times 512$ input.
$B$ denotes batch size.  Parameter counts exclude batch-normalisation
running statistics.}
\label{tab:architecture}
\small
\begin{tabular}{@{}llll@{}}
\toprule
\textbf{Module} & \textbf{Output shape} & \textbf{Key layers} & \textbf{Params} \\
\midrule
\multicolumn{4}{@{}l}{\textit{Dual Encoders (shared architecture, independent weights)}} \\
IR Encoder (ResNet-18 stages 1--4) & $16{\times}16{\times}512$ & conv, BN, ReLU, residual & 11.2\,M \\
RGB Encoder (ResNet-18 stages 1--4) & $16{\times}16{\times}512$ & (same) & 11.2\,M \\
IR FPN & $\{128{\times}128,\,64^2,\,32^2,\,16^2\}{\times}128$ & $1{\times}1$ conv + upsample & 0.8\,M \\
RGB FPN & (same) & (same) & 0.8\,M \\
\midrule
\multicolumn{4}{@{}l}{\textit{Feature Projection + CBAM Attention}} \\
Per-level $1{\times}1$ conv to $D{=}128$ & $H_l{\times}W_l{\times}128$ & conv $1{\times}1$ & 0.3\,M \\
CBAM (per level, per modality) & (same) & MLP + $7{\times}7$ conv & 0.1\,M \\
\midrule
\multicolumn{4}{@{}l}{\textit{Correlation Module (per pyramid level)}} \\
Dot-product + pyramid pooling & $H_l{\times}W_l{\times}H_l{\times}W_l$ & matmul, avg-pool & 0 \\
\midrule
\multicolumn{4}{@{}l}{\textit{Context Encoder}} \\
ResNet-18 (stages 1--3) on $I_{\mathrm{IR}}$ & $64{\times}64{\times}256$ & conv, BN, ReLU & 5.6\,M \\
Context projection & $64{\times}64{\times}128$ & conv $1{\times}1$ & 0.03\,M \\
\midrule
\multicolumn{4}{@{}l}{\textit{Iterative GRU Update Operator}} \\
Correlation look-up (radius $r{=}4$) & $H_l{\times}W_l{\times}(4{\cdot}81)$ & indexing & 0 \\
Displacement encoder & $H_l{\times}W_l{\times}128$ & 2 conv layers & 0.3\,M \\
Convolutional GRU ($3{\times}3$ gates) & $H_l{\times}W_l{\times}128$ & $z$/$r$/$\tilde{h}$ gates & 0.6\,M \\
Displacement head & $H_l{\times}W_l{\times}2$ & 2 conv layers & 0.03\,M \\
Confidence head & $H_l{\times}W_l{\times}1$ & 2 conv layers + sigmoid & 0.02\,M \\
\midrule
\multicolumn{4}{@{}l}{\textit{Upsampling Module}} \\
Convex upsampling (RAFT-style) & $H{\times}W{\times}2$ & $3{\times}3$ conv, softmax & 0.5\,M \\
\midrule
\textbf{Total} & & & $\sim$\textbf{31.5\,M} \\
\bottomrule
\end{tabular}
\end{table}

\paragraph{Encoder detail.}
Each ResNet-18 backbone follows the standard four-stage layout:
\texttt{conv1} ($7 \times 7$, stride 2) $\rightarrow$ \texttt{bn1}
$\rightarrow$ \texttt{relu} $\rightarrow$ \texttt{maxpool} (stride 2)
$\rightarrow$ \texttt{layer1} ($2 \times$ BasicBlock, 64 channels)
$\rightarrow$ \texttt{layer2} ($2 \times$ BasicBlock, 128 ch., stride 2)
$\rightarrow$ \texttt{layer3} ($2 \times$ BasicBlock, 256 ch., stride 2)
$\rightarrow$ \texttt{layer4} ($2 \times$ BasicBlock, 512 ch., stride 2).
Weights are initialised from \texttt{torchvision.models.resnet18
(weights=ResNet18\_Weights.DEFAULT)}.  For the IR encoder, the first
convolution is replaced by a $7 \times 7$ kernel with $C_{\text{in}}{=}1$.

\paragraph{Correlation look-up.}
At each GRU iteration the current displacement estimate
$\hat{\mathbf{d}}_t$ is used to index the correlation pyramid within a
local radius $r{=}4$, yielding $(2r+1)^2 = 81$ values per pyramid level.
With four pyramid levels the concatenated look-up vector has
$4 \times 81 = 324$ channels per pixel.

\paragraph{Convex upsampling.}
The displacement field predicted at $1/8$ resolution is upsampled to full
resolution using the \emph{convex upsampling} mechanism of RAFT
\cite{teed2020raft}: for each output pixel the network predicts a
$8 \times 8$ weight mask over its coarse-resolution neighbours via a
$3 \times 3$ convolution followed by softmax, and the final displacement
is a weighted combination.  This avoids the checkerboard artefacts of
transposed convolutions and the blurring of bilinear interpolation.

% ========================================================================
% FIGURE PLACEHOLDER
% ========================================================================
\begin{figure}[t]
  \centering
  % -----------------------------------------------------------------
  % PLACEHOLDER --- replace with actual TikZ or included graphic.
  %
  % Description of the figure to be created:
  %
  % The diagram flows left to right.  On the far left, two input images
  % are stacked vertically: "I_IR (H x W x 1)" (top) and "I_RGB
  % (H x W x 3)" (bottom).  Arrows lead from each image into its own
  % encoder box (light blue for IR, light orange for RGB), each labelled
  % "ResNet-18 + FPN".  Each encoder outputs a feature pyramid shown as
  % four rectangles of decreasing size.
  %
  % The two pyramids feed into a central "Correlation Volume" block
  % (green), from which arrows flow into the "Iterative GRU Refiner"
  % block (purple) that also receives input from a "Context Encoder"
  % box attached to I_IR.  Inside the GRU block, a circular arrow
  % indicates recurrence (T iterations).
  %
  % The GRU block has two output arrows: one to a "Displacement Head
  % (H x W x 2)" box and one to a "Confidence Head (H x W x 1)" box,
  % both coloured yellow.  The displacement arrow continues into a
  % "Differentiable Warp (grid_sample)" box (red) which also receives
  % I_RGB and outputs "I_RGB^warp (H x W x 3)".
  %
  % A dashed box surrounds the GRU + heads and is labelled "Coarse-to-
  % Fine: repeat at each pyramid level".  A small legend in the corner
  % maps colours to module types.
  % -----------------------------------------------------------------
  \fbox{\parbox{0.92\textwidth}{\centering\vspace{4em}%
    \textbf{[Figure placeholder -- see source comments for detailed
    description of the architecture diagram to be inserted here.]}%
    \vspace{4em}}}
  \caption{Overview of IRAlign-Net.  Dual encoders extract modality-specific
  feature pyramids from the IR and RGB inputs.  A 4D correlation volume
  enables dense matching, and a convolutional GRU iteratively refines
  the displacement field in a coarse-to-fine manner.  Parallel output
  heads predict the $H \times W \times 2$ displacement field and the
  $H \times W \times 1$ confidence mask.  A differentiable warping
  layer applies the predicted displacement to the RGB image.}
  \label{fig:architecture}
\end{figure}

% ========================================================================
\section{Implementation Considerations}
\label{sec:implementation}

% ------------------------------------------------------------------------
\subsection{Differentiable Warping in PyTorch}
\label{subsec:grid_sample}

The warping operation of Equation~\eqref{eq:warp} is implemented with a
two-line pattern in PyTorch:

\begin{verbatim}
  # displacement: (B, 2, H, W) in pixel units
  # Build normalised sampling grid
  grid_y, grid_x = torch.meshgrid(
      torch.linspace(-1, 1, H, device=disp.device),
      torch.linspace(-1, 1, W, device=disp.device),
      indexing='ij')
  base_grid = torch.stack([grid_x, grid_y], dim=-1)  # (H,W,2)
  # Convert pixel displacement to normalised displacement
  norm_disp = torch.zeros_like(displacement)
  norm_disp[:, 0] = displacement[:, 0] / (W / 2)
  norm_disp[:, 1] = displacement[:, 1] / (H / 2)
  sample_grid = base_grid + norm_disp.permute(0,2,3,1)
  warped = F.grid_sample(
      rgb, sample_grid, mode='bilinear',
      padding_mode='border', align_corners=True)
\end{verbatim}

Key implementation notes:
\begin{itemize}
  \item \texttt{grid\_sample} expects the grid in $(x,y)$ order with
        $x \in [-1,1]$ mapping to the width axis.  A common source of
        bugs is swapping $x$ and $y$.
  \item \texttt{padding\_mode='border'} replicates edge pixels for
        out-of-bounds samples, which is preferable to zero-padding for
        registration tasks because it avoids black borders that
        contaminate downstream losses.
  \item \texttt{align\_corners=True} ensures that pixel centres at the
        image boundary map to $\pm 1$ rather than the half-pixel-shifted
        convention, which matters when displacement magnitudes are large.
  \item Gradients flow through both the grid (enabling displacement
        learning) and the input image (enabling joint fine-tuning of
        upstream feature extractors), courtesy of the bilinear
        interpolation kernel's differentiability \cite{jaderberg2015spatial}.
\end{itemize}

% ------------------------------------------------------------------------
\subsection{Memory Management for Correlation Volumes}
\label{subsec:memory}

The all-pairs correlation volume $\mathbf{C} \in \mathbb{R}^{H' \times
W' \times H' \times W'}$ is the primary memory bottleneck.  For
$H'{=}W'{=}64$ the volume occupies 67\,MB at FP32.  Scaling to
$H'{=}W'{=}128$ ($1/4$ resolution for a $512$ input) would require
$128^4 \times 4 \approx 1.07$\,GB, which is prohibitive for multi-scale
operation on consumer GPUs.

We adopt three complementary strategies:
\begin{enumerate}
  \item \textbf{Compute at $1/8$ resolution.}  The correlation volume is
        built only at the $1/8$ level; finer levels use the upsampled
        displacement from the coarser level as initialisation and perform
        local correlation within a search window of radius $r{=}4$,
        reducing memory from $O(H'^2 W'^2)$ to $O(H' W' (2r+1)^2)$.
  \item \textbf{FP16 mixed precision.}  The correlation volume is
        stored in half precision (33.5\,MB at $64^2$), and the GRU
        update operator runs in FP16 via PyTorch's
        \texttt{torch.cuda.amp.autocast} \cite{micikevicius2018mixed}.
        The displacement head output is cast back to FP32 before the
        warping step to preserve sub-pixel accuracy.
  \item \textbf{Gradient checkpointing.}  During training, intermediate
        activations within the GRU unrolling are recomputed rather than
        stored, trading a $\sim$25\% increase in compute time for a
        $\sim$60\% reduction in activation memory.  This is implemented
        via \texttt{torch.utils.checkpoint.checkpoint}.
\end{enumerate}

With these optimisations, the full IRAlign-Net trains on two $512 \times
512$ image pairs per GPU at batch size 8 (4 GPUs, gradient accumulation 1)
on NVIDIA A100 (40\,GB) GPUs, and at batch size 4 on a single RTX 3090
(24\,GB).

% ------------------------------------------------------------------------
\subsection{Inference Speed for Real-Time Video}
\label{subsec:speed}

For deployment on a synchronised IR--RGB video feed, latency and
throughput are critical.  Table~\ref{tab:speed} summarises measured
timings.

\begin{table}[t]
\centering
\caption{Inference latency on a single NVIDIA RTX 3090 GPU
($512 \times 512$ input, FP16, batch size 1).}
\label{tab:speed}
\small
\begin{tabular}{@{}lcc@{}}
\toprule
\textbf{Component} & \textbf{Time (ms)} & \textbf{Fraction} \\
\midrule
Dual encoders + FPN       &  8.2 & 28\% \\
Correlation volume        &  3.1 & 11\% \\
GRU refinement (12 iter.) & 14.6 & 50\% \\
Convex upsampling         &  1.5 &  5\% \\
Confidence head           &  0.4 &  1\% \\
Grid sample warp          &  1.4 &  5\% \\
\midrule
\textbf{Total}            & \textbf{29.2} & 100\% \\
\bottomrule
\end{tabular}
\end{table}

At 29.2\,ms per frame the network achieves approximately 34\,FPS, which
is adequate for 30\,Hz video.  If faster inference is required, two
strategies are effective:
\begin{itemize}
  \item \textbf{Reducing GRU iterations.}  Cutting from 12 to 6
        iterations reduces latency to $\sim$21\,ms (47\,FPS) at the cost
        of approximately 0.3\,px increase in endpoint error.
  \item \textbf{TensorRT export.}  Converting the PyTorch model to
        TensorRT with INT8 quantisation yields a further $2\times$
        speed-up on Ampere-architecture GPUs, reaching approximately
        80\,FPS.  Quantisation-aware training \cite{jacob2018quantization}
        is recommended to maintain accuracy.
  \item \textbf{Temporal warm-starting.}  In video mode, the previous
        frame's displacement field initialises the current frame's GRU
        iterations, reducing the number of iterations needed for
        convergence from 12 to 3--4 in regions with smooth motion.
\end{itemize}

% ========================================================================
\section{Design Alternatives and Justification}
\label{sec:alternatives}

Several architectural alternatives were considered during the design of
IRAlign-Net.

\paragraph{Single-encoder U-Net.}
Concatenating the four-channel input $[I_{\mathrm{IR}}; I_{\mathrm{RGB}}]$
and using a standard U-Net encoder--decoder \cite{ronneberger2015unet}
reduces the parameter count to $\sim$18\,M and improves throughput to
$\sim$55$\,FPS.  However, the shared early layers are forced to reconcile
the disparate IR and RGB statistics, leading to weaker feature
representations and a 0.4--0.7\,px increase in endpoint error on our
benchmark (see Chapter~\ref{ch:experiments}).

\paragraph{Direct regression without correlation volume.}
A purely feed-forward decoder that directly regresses displacements from
concatenated encoder features eliminates the cost-volume computation
entirely.  In our experiments this approach converges faster in early
epochs but plateaus at a higher error, consistent with findings in the
optical-flow literature \cite{sun2018pwcnet}.  The correlation volume
provides an explicit matching signal that dramatically improves
generalisation to displacement magnitudes not seen during training.

\paragraph{Transformer-based matching.}
Recent work in optical flow (FlowFormer \cite{huang2022flowformer},
GMFlow \cite{xu2022gmflow}) replaces the 4D correlation volume with
cross-attention between feature tokens.  While this offers global
receptive fields and handles large displacements elegantly, the quadratic
memory cost of full attention at $64^2{=}4096$ tokens is comparable to the
correlation volume, and our preliminary experiments showed no statistically
significant accuracy improvement on IR--RGB data with only a few thousand
training pairs.  We therefore retain the convolutional correlation
approach for its simplicity and well-understood training dynamics, but note
that transformer-based matching may become advantageous as larger IR--RGB
datasets become available.

% ========================================================================
\section{Summary}
\label{sec:arch_summary}

This chapter has presented IRAlign-Net, a dense alignment architecture
tailored to infrared--visible image registration.  The key design choices
are:
\begin{enumerate}
  \item \textbf{Dual modality-specific encoders} with Feature Pyramid
        Networks to respect the distinct statistical properties of IR and
        RGB imagery.
  \item \textbf{4D correlation volumes} with pyramid pooling to provide
        an explicit, fine-grained matching signal across modalities.
  \item \textbf{Iterative GRU-based refinement} in a coarse-to-fine
        cascade, enabling both large-displacement detection and sub-pixel
        precision.
  \item \textbf{Confidence mask prediction} that allows the network to
        self-assess reliability and prevents gradient contamination from
        inherently ambiguous regions.
  \item \textbf{Semantic-mask gating and spatial attention} to focus
        alignment effort on safety-critical objects and informative
        image regions.
\end{enumerate}

\noindent
The resulting network contains approximately 31.5\,M parameters, runs at
$\sim$34\,FPS on a single GPU, and---as demonstrated in
Chapter~\ref{ch:experiments}---achieves state-of-the-art pixel-wise
alignment accuracy on our IR--RGB benchmark.
